% ID: SEL-D0032-OA-E03
% ESTADO: cotejada. Segunda pasada visual y borrador compilado revisados.
% Fuente: Selectividad/matematicas_aplicadas/titular_junio_EXAMEN 5-ANDALUCÍA.pdf
% SHA-256: 5f24b200fb6cf34f6d72016ec3dc35787b4074de9b7752677e5b0d959ed0dfcd
% Página: 1; opción A; ejercicio original 3.

Se sabe que el 90\% de los alumnos de un centro docente está interesado por las redes sociales, el 60\% está interesado por sus notas y el 55\% por ambas cuestiones. Se elige al azar un alumno de ese centro.
\begin{enumerate}[label=\alph*)]
\item \textbf{(1 punto)} \textquestiondown{}Cuál es la probabilidad de que dicho alumno esté interesado por alguna de las dos cuestiones?
\item \textbf{(1 punto)} Calcule la probabilidad de que esté interesado por sus notas, sabiendo que no está interesado por las redes sociales.
\item \textbf{(0.5 puntos)} Calcule la probabilidad de que no esté interesado por ninguna de estas dos cuestiones.
\end{enumerate}
